% Generated by Listas/make_statements.py from pratica-lecture-07.md, item 23. Do not edit.
\textbf{Source.} Tadelis (2013), \emph{Game Theory: An Introduction}, Exercise 13.5 ``Mixed Strategies'', p.~286 {[}PDF p.~305{]}. Reproduced verbatim. \textbf{Answer in the source:} no (none in the book; Princeton UP's website has about half, not checked).



\begin{quote}

\textbf{Mixed Strategies:} Consider a seller who offers a single private value good using a first-price sealed-bid auction. There are two potential buyers, each with a valuation that can take on one of two values, \(\theta_i \in \{4, 8\}\), each value occurring with an equal probability of \(\frac{1}{2}\). The players' values are independently drawn. Call the type with \(\theta_i = 4\) the ``low'' type and \(\theta_i = 8\) the ``high'' type.



\begin{enumerate}

\def\labelenumi{\alph{enumi}.}

\item

  Show that in a Bayesian Nash equilibrium the low must be using a pure strategy bid \(b_L = 4\). (Hint: First show that in a Bayesian Nash equilibrium the low type cannot be using a pure strategy bid \(b_L < 4\). Next show that in a Bayesian Nash equilibrium the low type cannot be using a mixed strategy with two different bids.)

\item

  Show that given the fact described in (a), in any Bayesian Nash equilibrium the high type will never choose a bid \(b_H > 6\).

\item

  Show that the low type choosing \(b_L = 4\) and the high type choosing \(b_H\) uniformly from the interval \([4, 6]\) together constitute a Bayesian Nash equilibrium.

\item

  Calculate the expected revenue to the seller given the Bayesian Nash equilibrium in (c). How does it compare to the expected revenue he would obtain from an English auction?

\end{enumerate}

\end{quote}



\textbf{Note on the book (item 23).} In part (c) the uniform distribution on \([4,6]\) is not an equilibrium strategy of

the high type (the payoff \((8-b)(b-2)/4\) is not constant on \([4,6]\)); the indifference condition gives the cdf

\((b-4)/(8-b)\) on \([4,6]\), whose revenue equals the English auction's \(5\). See the solutions PDF, item 23.

